% =============================================================================
% glossary/entries.tex: abbreviation and glossary entries.
% Loaded at the end of preamble/glossary.tex (see that file for all the commands you can use in the text).
%
% Only entries used in the text appear in Appendix A, sorted alphabetically, so defining an entry you never use is harmless. Keep this file grouped by kind and alphabetical by label so entries are easy to find.
%
% THE THREE KINDS OF ENTRY
%
% 1. Abbreviation only -> A.1 Abbreviations
%      \newabbreviation{label}{SHORT}{long form}
%    Write the long form in lower case unless it is a proper noun; \Gls capitalises it at the start of a sentence.
%      \newabbreviation[longplural={...}, shortplural={...}]{label}{SHORT}{long form}
%
% 2. Defined term only -> A.2 Glossary
%      \newglossaryentry{label}{name={term}, description={definition}}
%    Optional keys: plural={...} for an irregular plural; sort={...} to sort under a different spelling.
%
% 3. Abbreviation that also has a definition -> both lists, linked
%      \newtermabbr{label}{SHORT}{long form}{definition}
%    In the text, \gls{label} gives the abbreviation and \gls{def-label} the defined term. The A.1 line ends "(see definition)" and links to A.2.
%
% Labels: lower case, no spaces; hyphens are fine (e.g. rug-pull). The label is only used in \gls{...}, never printed. Commas or "=" inside a value need braces: description={a, b, and c}.
% =============================================================================

% --- 1. Abbreviations ----------------------------------------------------------------------


% --- 3. Abbreviations with a definition ------------------------------------------------------
% The fourth argument (the definition) may start on the next line, but there must be no blank line between it and the third argument.


% --- 2. Defined terms ------------------------------------------------------------------------

